
\section{射影空间中的\BookRecordIndex{T0386}有理曲线}\label{sec:7-4}
\zhindex{sheyingkongjianzhongdeyouliquxian}{射影空间中的有理曲线}

射影空间中的曲线具有错综而迷人的性质，对它们的研究也有悠久历史。不过我们将看到，
若目的仅仅是计算 $\mathbb{CP}^n$ 中的亏格零\BookRecordIndex{T0143}格罗莫夫--威滕不变量，那么它们的大多数
复杂行为都无关紧要。

$n=1$ 的情形并不有趣。这是因为，当 $M=\mathbb{CP}^1$ 时，带 $k$ 个\BookRecordIndex{T0066}标记点的次数
$d>1$ 全纯球面（即有理映射）\BookRecordIndex{T0246}模空间维数为
\[
  \mu(dL,k)=2(k+2d-2)>2k,
\]
所以除非 $d=1$，\BookRecordIndex{T0145}格罗莫夫--威滕伪循环
\[
  \operatorname{ev}:\mathcal M^*_{0,k}(dL;J)\longrightarrow(\mathbb{CP}^1)^k
\]
都是平凡的。不过，计数更高亏格的曲线确实会得到有趣的不变量。它们称为
\BookRecordIndex{T0045}Riemann--Hurwitz 数，并与 \BookRecordIndex{T0050}Toda 晶格有深刻联系；见 Okounkov--Pandharipande
\srcite{311, 312}。

以下假设 $n\ge2$。这时，带 $k$ 个\BookRecordIndex{T0066}标记点的次数 $d$ 全纯球面\BookRecordIndex{T0246}模空间的虚维数为
\[
  \dim\mathcal M^*_{0,k}(dL;J_0)=2m,
  \qquad
  m:=d(n+1)+n+k-3.
\]
下面将看到，这正是其实际维数，而且 $J_0$ 在定义~\ref{def:6-2-1} 的意义下正则。为了
得到非零不变量，必须假设 $m\le nk$，等价地，
\[
  (n-1)k\ge d(n+1)+n-3.
\]
例~\ref{ex:7-4-11} 表明，即使满足维数条件，不变量仍可能为零。不过我们将看到，当
$k=3d-1$ 时它总是非零。本节末尾还会讨论其他非零情形。先证明下面的定理。

\begin{theorem}\label{thm:7-4-1}
设 $n\ge2$ 且 $d,k\ge1$。则\BookRecordIndex{T0283}求值映射
\[
  \operatorname{ev}_{J_0}:\mathcal M^*_{0,k}(dL;J_0)
  \longrightarrow(\mathbb{CP}^n)^k
\]
是实维数 $2d(n+1)+2n+2k-6$ 的\BookRecordIndex{T0337}伪循环。此外，若 $k=3d-1$，它代表非零同调类。
\end{theorem}

这个定理的两个部分有不同的证明。第一个断言是初等的；它成立是因为复流形
$(\mathbb{CP}^n,J_0)$ 有一个传递作用的复自同构群
$\operatorname{PSL}(n+1)$。第二个断言较难证明。它很容易由 $n=2$ 的情形\BookRecordIndex{T0326}推出，而后一
情形来自 Kontsevich 关于 $\mathbb{CP}^2$ 中次数 $d$ 球面数目的递推公式
（命题~\ref{prop:7-5-11}）。但若要建立这个二维结论，就必须使用第~10 章的\BookRecordIndex{T0256}粘合定理，
或代数几何中一些颇为精细的方法。本节给出的定理~\ref{thm:7-4-1} 的证明是完整的，不过
确实需要定理~\ref{thm:2-6-4} 的\BookRecordIndex{T0055}伴随公式，具体采用如下形式。

\begin{remark}\label{rem:7-4-2}
设 $u:\Sigma\to M^4$ 是单纯 \BookRecordIndex{T0003}$J$-全纯曲线。回忆，$\Sigma$ 中两个不同点
$z_1,z_2$ 所成的无序对称为 $u$ 的一个\emph{\BookRecordIndex{T0167}横截自交点}，如果
\[
  u(z_1)=u(z_2)=:x,
  \qquad
  T_xM=\operatorname{im}du(z_1)\oplus\operatorname{im}du(z_2).
\]
若 $u:\Sigma\to\mathbb{CP}^2$ 是一条亏格为 $g$、次数为 $d$、有 $\delta$ 个\BookRecordIndex{T0167}横截自交点
的单纯 $J_0$-全纯曲线，则由定理~\ref{thm:2-6-4}，
\[
  g+\delta\le\frac{(d-1)(d-2)}2,
\]
并且等号成立当且仅当 $u$ 是浸入，且只有\BookRecordIndex{T0167}横截自交点。这里允许多重交点；也就是说，可能
存在 $k$ 个点 $z_1,\ldots,z_k$，使每个无序对 $z_i,z_j$（$i\ne j$）都满足上述条件。
\end{remark}

\paragraph{标准复结构的\BookRecordIndex{T0408}正则性。}
我们将证明 $J_0\in\mathcal J_{\mathrm{reg}}(\mathbb{CP}^n,\omega_0)$，其中 $\omega_0$
表示 Fubini--Study 度量的\BookRecordIndex{T0363}辛形式。这个证明也适用于自同构群 $G$ 传递作用的任意 \BookRecordIndex{T0207}凯勒
流形 $(M,J_0,\omega_0)$，例如格拉斯曼流形和对称空间 $G/T$。当
$M=\mathbb{CP}^n$ 时，相应的群是 $G=\operatorname{PSL}(n+1)$。回忆，
$J_0\in\mathcal J_{\mathrm{reg}}(M,\omega_0)$ 当且仅当：对每个单纯 $J_0$-全纯球面
$u:S^2\to M$，算子 $D_u$ 都满射；而且对每棵\BookRecordIndex{T0084}带标号树 $T$ 和每族满足稳定性条件
\eqref{eq:6-1-1} 的球面同调类 $\{A_\alpha\}_{\alpha\in T}$，\BookRecordIndex{T0065}边求值映射
\[
  \operatorname{ev}^E:
  \mathcal M^*(\{A_\alpha\};J_0)\times Z(T)
  \longrightarrow M^E
\]
都横截于 $\Delta^E$。下面的证明还表明，在我们的假设下，连多重覆盖的
$J_0$-全纯球面也是正则的。

\begin{proposition}\label{prop:7-4-3}
设 $(M,J_0,\omega_0)$ 是紧 \BookRecordIndex{T0207}凯勒 流形，Lie 群 $G$ 通过全纯微分同胚传递作用于 $M$。
则
\[
  J_0\in\mathcal J_{\mathrm{reg}}(M,\omega_0),
\]
而且 $M$ 中每个非常值 $J_0$-全纯球面的 Chern 数都至少为 $2$。
\end{proposition}

\begin{proof}
证明分四步。

\emph{步骤 1. 对每个 $x\in M$，无穷小作用 $\mathfrak g\to T_xM$ 都满射。}

这是因为，$TM$ 中由 $\mathfrak g$ 张成的子丛是一个\BookRecordIndex{T0123}分布，其积分子流形就是 $G$ 的轨道。
根据假设，每个轨道都是整个流形 $M$。

\emph{步骤 2. 每个 $J_0$-全纯球面 $u:S^2\to M$ 都正则。}

由步骤 1，切丛 $TM$ 在每一点都由全局全纯向量场
\[
  X_\xi(x):=\left.\frac d{dt}\right|_{t=0}\exp(t\xi)x,
  \qquad \xi\in\mathfrak g,
\]
张成。现在设 $u:S^2\to M$ 是 $J_0$-全纯球面。由 Grothendieck 定理，拉回切丛分裂成
全纯线丛的直和
\[
  u^*TM=L_1\oplus\cdots\oplus L_n.
\]
记 $\pi_i:u^*TM\to L_i$ 为投影。每个全纯向量场
$X\in\operatorname{Vect}(M)$ 都下降为 $L_i$ 的全纯截面
$s_i:=\pi_iX(u)$。因此，由步骤 1，每个子丛 $L_i$ 都有非零全纯截面，从而次数非负。于是
由引理~\ref{lem:3-3-2}，$D_u$ 满射。

\emph{步骤 3. 每个非常值 $J_0$-全纯球面的 Chern 数都至少为 $2$。}

设 $A\in H_2(M;\mathbb Z)$ 是可由非常值 $J_0$-全纯球面代表的同调类。带一个\BookRecordIndex{T0066}标记点的
\BookRecordIndex{T0246}模空间 $\mathcal M_{0,1}(A;J_0)$ 非空，而且由步骤 2，它正则，维数为
$2n+2c_1(A)-4$。由于 $G$ 的作用传递，\BookRecordIndex{T0283}求值映射
\[
  \operatorname{ev}:\mathcal M_{0,1}(A;J_0)\longrightarrow M
\]
满射。\sourcefootnote{原书此处把\BookRecordIndex{T0283}求值映射的定义域印作
$\mathcal M^*_{0,1}(A;J_0)$；但上一句讨论的是包含给定非常值球面的完整\BookRecordIndex{T0246}模空间
$\mathcal M_{0,1}(A;J_0)$，命题结论也包括多重覆盖球面。故据此补去星号。}
所以这个\BookRecordIndex{T0246}模空间的维数至少为 $2n$，从而 $c_1(A)\ge2$。

\emph{步骤 4. 对每棵\BookRecordIndex{T0084}带标号树 $T$ 和每族球面同调类
$\{A_\alpha\}_{\alpha\in T}\subset H_2(M;\mathbb Z)$，\BookRecordIndex{T0065}边求值映射}
\[
  \operatorname{ev}^E:
  \mathcal M(\{A_\alpha\};J_0)\times Z(T)\longrightarrow M^E
\]
\emph{都横截于 $\Delta^E$。}

这是步骤 1 的直接推论。以下论证给出细节。给定
\[
  \mathbf u=\{u_\alpha\}_{\alpha\in T}
  \in\mathcal M(\{A_\alpha\};J_0),
  \qquad
  \mathbf z=
  \bigl(\{z_{\alpha\beta}\}_{\alpha E\beta},
  \{z_i\}_{1\le i\le k}\bigr)\in Z(T),
\]
记
\[
  \mathbf x:=\operatorname{ev}^E(\mathbf u,\mathbf z)\in\Delta^E.
\]
这意味着 $\mathbf x=\{x_{\alpha\beta}\}_{\alpha E\beta}$，其中
$x_{\alpha\beta}:=u_\alpha(z_{\alpha\beta})$，且对 $\alpha E\beta$ 有
$x_{\alpha\beta}=x_{\beta\alpha}$。记
\[
  V_{\mathbf x}:=
  \bigl\{\{X_{\xi_\alpha}(x_{\alpha\beta})\}_{\alpha E\beta}
  \mid \xi_\alpha\in\mathfrak g\bigr\}
  \subset T_{\mathbf x}M^E.
\]
我们将证明
\begin{equation}\label{eq:7-4-1}
  V_{\mathbf x}+T_{\mathbf x}\Delta^E=T_{\mathbf x}M^E.
\end{equation}
由于群 $G^T$ 作用于\BookRecordIndex{T0246}模空间 $\mathcal M(\{A_\alpha\};J_0)$，子空间
$V_{\mathbf x}$ 包含在微分 $d\operatorname{ev}^E(\mathbf u,\mathbf z)$ 的像中。因此
式~\eqref{eq:7-4-1} 表明 $\operatorname{ev}^E$ 横截于 $\Delta^E$，正如所断言的。

$T_{\mathbf x}\Delta^E$ 的一个补空间由满足
\[
  \alpha E\beta\qquad\Longrightarrow\qquad
  v_{\alpha\beta}+v_{\beta\alpha}=0
\]
的向量
$\mathbf v=\{v_{\alpha\beta}\}_{\alpha E\beta}\in T_{\mathbf x}M^E$ 张成。因此，为证明
式~\eqref{eq:7-4-1}，只须证明：对每个这样的 $\mathbf v$，都存在
$\{\xi_\alpha\}_{\alpha\in T}\in\mathfrak g^T$，使得对 $\alpha E\beta$ 有
\begin{equation}\label{eq:7-4-2}
  v_{\alpha\beta}
  =X_{\xi_\alpha}(x_{\alpha\beta})
   -X_{\xi_\beta}(x_{\beta\alpha}).
\end{equation}
我们对边数作归纳。若没有边，便无须证明。因而假设 $e(T)>0$，并假定对每棵边数至多为
$e(T)-1$ 的\BookRecordIndex{T0417}子树 $T_0\subset T$，式~\eqref{eq:7-4-2} 已经成立。取树的一个端点
$\alpha_0\in T$，并令 $\beta_0\in T_0:=T\setminus\{\alpha_0\}$ 是满足
$\alpha_0E\beta_0$ 的唯一顶点。归纳假设表明，存在
$\{\xi_\alpha\}_{\alpha\in T_0}\in\mathfrak g^{T_0}$，使式~\eqref{eq:7-4-2}
对所有满足 $\alpha E\beta$ 的 $\alpha,\beta\in T_0$ 成立。由步骤 1，存在
$\xi_{\alpha_0}\in\mathfrak g$，使式~\eqref{eq:7-4-2} 对
$\alpha=\alpha_0,\beta=\beta_0$ 成立。由于只有一条边邻接 $\alpha_0$，该式于是对所有
$\alpha E\beta$ 成立。这证明了式~\eqref{eq:7-4-1}。
\end{proof}

\begin{remark}\label{rem:7-4-4}
在 $M=\mathbb{CP}^n$ 且 $G=\operatorname{PSL}(n+1)$ 的情形，命题~\ref{prop:7-4-3}
证明中的全纯向量场 $X_\xi$ 有如下显式形式。把元素
$\ell\in\mathbb{CP}^n$ 看成 $\mathbb C^{n+1}$ 中的一条直线，并把切空间等同于
\[
  T_\ell\mathbb{CP}^n
  =\operatorname{Hom}(\ell,\mathbb C^{n+1}/\ell).
\]
那么，Lie 代数
$\mathfrak g:=\operatorname{Lie}(G)=\mathfrak{sl}(n+1)$ 的无穷小作用为
\[
  X_\xi(\ell)=\pi_\ell\circ\xi\circ\iota_\ell,
  \qquad \xi\in\mathfrak g,
\]
其中 $\iota_\ell:\ell\to\mathbb C^{n+1}$ 是包含映射，
$\pi_\ell:\mathbb C^{n+1}\to\mathbb C^{n+1}/\ell$ 是投影。这些向量场是切丛的全纯截面，
而且显然在每个 $\ell$ 处张成切空间 $T_\ell\mathbb{CP}^n$。
\end{remark}

下一个结果给出\BookRecordIndex{T0143}格罗莫夫--威滕不变量非零的一个充分条件。其思路是：若上同调类 $a_i$ 是
复子流形 $X_i$ 的 Poincaré 对偶，且乘积 $X_1\times\cdots\times X_k$ 横截于所有\BookRecordIndex{T0283}求值映射
\[
  \operatorname{ev}_{T,J_0}:\mathcal M^*_{0,T}(A;J_0)\longrightarrow M^k,
\]
那么每条穿过各 $X_i$ 的 $J_0$-全纯球面对\BookRecordIndex{T0143}格罗莫夫--威滕不变量的贡献都是正的。因此，只须
找到一个解。

\begin{proposition}\label{prop:7-4-5}
设 $(M,J_0,\omega_0)$ 和 $G$ 如命题~\ref{prop:7-4-3}，并假设
$A\in H_2(M;\mathbb Z)$ 与 $a_1,\ldots,a_k\in H^*(M)$ 满足维数条件
\eqref{eq:7-1-2}。进一步假设存在紧复子流形 $X_1,\ldots,X_k\subset M$、
$J_0$-全纯球面 $u:\mathbb{CP}^1\to M$，以及 $\mathbb{CP}^1$ 中两两不同的点所成的元组
$\mathbf z=(z_1,\ldots,z_k)$，使得对 $i=1,\ldots,k$ 有
\[
  a_i=\operatorname{PD}([X_i]),
  \qquad
  u(z_i)\in X_i,
\]
且\BookRecordIndex{T0283}求值映射
\[
  \operatorname{ev}_{J_0}:\mathcal M^*_{0,k}(A;J_0)\longrightarrow M^k
\]
在 $[u,\mathbf z]$ 处横截于乘积 $X_1\times\cdots\times X_k$。则
\[
  \mathrm{GW}^M_{A,k}(a_1,\ldots,a_k)>0.
\]
\end{proposition}

\begin{proof}
先假设乘积 $\mathbf X:=X_1\times\cdots\times X_k$ 对每棵 $k$-\BookRecordIndex{T0084}带标号树 $T$ 和每个球面
同调类 $B\in H_2(M;\mathbb Z)$，都横截于\BookRecordIndex{T0283}求值映射
\[
  \operatorname{ev}_{T,J_0}:\mathcal M^*_{0,T}(B;J_0)
  \longrightarrow\mathbf M:=M^k.
\]
那么集合
\[
  \mathcal M^*_{0,k}(A,\mathbf X;J_0)
  :=\operatorname{ev}_{J_0}^{-1}(\mathbf X)
  =\bigl\{[u,\mathbf z]\in\mathcal M^*_{0,k}(A;J_0)\mid
  u(z_i)\in X_i\bigr\}
\]
是有限的；按定义，其元素带适当符号的计数等于
$\mathrm{GW}^M_{A,k}(a_1,\ldots,a_k)$。由假设，
$\mathcal M^*_{0,k}(A,\mathbf X;J_0)$ 非空。此外，因为 $J_0$ 可积，且
$\mathbf X$ 是 $\mathbf M$ 的复子流形，所以对每个
$[u,\mathbf z]\in\mathcal M^*_{0,k}(A,\mathbf X;J_0)$，在
$\mathbf x:=(u(z_1),\ldots,u(z_k))\in\mathbf X$ 处，直和分解
\[
  T_{\mathbf x}\mathbf M
  =T_{\mathbf x}\mathbf X
   \oplus\operatorname{im}d\operatorname{ev}_{J_0}([u,\mathbf z])
\]
中的每个空间都是 $T_{\mathbf x}\mathbf M$ 的复子空间。因此，每个解
$[u,\mathbf z]$ 对不变量贡献 $1$，从而
\[
  \mathrm{GW}^M_{A,k}(a_1,\ldots,a_k)
  =\#\mathcal M^*_{0,k}(A,\mathbf X;J_0)>0.
\]

现在考虑任意子流形 $\mathbf X$，并令
$\mathbf G_{\mathrm{reg}}\subset\mathbf G:=G^k$ 表示所有满足如下条件的元组
$\mathbf g=(g_1,\ldots,g_k)\in\mathbf G$ 所成的集合：对每个 $T$ 和 $B$，
$\mathbf g\mathbf X$ 都横截于 $\operatorname{ev}_{T,J_0}$。我们证明
$\mathbf G_{\mathrm{reg}}$ 在 $\mathbf G$ 中是剩余的。更一般地，以定义在光滑流形 $U$ 上的
任意光滑映射 $f:U\to\mathbf M$ 代替 $\operatorname{ev}_{T,J_0}$，并构造空间
\[
  Y_f:=\bigl\{(\mathbf g,\mathbf x,u)\in
  \mathbf G\times\mathbf X\times U\mid f(u)=\mathbf g\mathbf x\bigr\}.
\]
这是 $\mathbf G\times\mathbf X\times U$ 的光滑子流形。事实上，由于 $\mathbf G$ 传递作用于
$\mathbf M$，映射
\[
  \mathbf G\times\mathbf X\times U\longrightarrow\mathbf M\times\mathbf M,
  \qquad
  (\mathbf g,\mathbf x,u)\longmapsto(\mathbf g\mathbf x,f(u))
\]
横截于对角线。元素 $\mathbf g\in\mathbf G$ 是投影 $Y_f\to\mathbf G$ 的\BookRecordIndex{T0409}正则值，当且仅当
$\mathbf g\mathbf X$ 横截于 $f$。所以，由 \BookRecordIndex{T0296}萨德定理，对每个光滑映射
$f:U\to\mathbf M$，所有使 $\mathbf g\mathbf X$ 横截于 $f$ 的
$\mathbf g\in\mathbf G$ 所成的集合 $\mathbf G_{\mathrm{reg}}(f)$ 都是剩余的。因而上述集合
$\mathbf G_{\mathrm{reg}}$ 在 $\mathbf G$ 中是剩余的，正如所断言的。

由我们的假设，只要 $\mathbf g\in\mathbf G_{\mathrm{reg}}$ 充分接近恒等元 $\mathbf1$，就有
$\mathcal M^*_{0,k}(A,\mathbf g\mathbf X;J_0)\ne\varnothing$。于是由证明的第一部分，
该不变量为正。这就证明了命题。
\end{proof}

\begin{remark}\label{rem:7-4-6}
考察射影特殊线性群
$G=\operatorname{PSL}(n+1,\mathbb C)$ 的作用，可以对 $\mathbb{CP}^n$ 中的\BookRecordIndex{T0386}有理曲线
获得一些有用认识。\sourcefootnote{原书把
$G=\operatorname{PSL}(n+1,\mathbb C)$ 称为“general linear group”；但该记号表示射影特殊
线性群，故据群的明确定义补正名称。}\BookRecordIndex{T0246}模空间
$\widetilde{\mathcal M}_{0,0}(\mathbb{CP}^n,dL;J_0)$ 中的每个元素 $u$ 都具有形式
\[
  u(z)=[u_0(z):\cdots:u_n(z)],
  \qquad z\in\mathbb C\cup\{\infty\},
\]
其中 $u_i$ 是次数不超过 $d$ 的多项式。因此，以 $u_i$ 的系数作为齐次坐标，可把
$\widetilde{\mathcal M}_{0,0}(\mathbb{CP}^n,dL;J_0)$ 等同于射影空间
$\mathbb{CP}^N$ 的一个开稠密子集，其中
\[
  N=(d+1)(n+1)-1.
\]
我们并没有得到整个 $\mathbb{CP}^N$，因为其中有些元素是退化的，并不代表正确的同调类；
例如当 $k<d$ 时的 $[1:z^k:0:\cdots:0]$。

对这样的 $u$，令 $F(u)$ 为一个 $(n+1)\times(d+1)$ 矩阵，其第 $i$ 行由多项式 $u_i$ 的
系数组成。这个矩阵在乘以一个标量的意义下良定义。记 $r(u)$ 为 $F(u)$ 的秩。于是
$r(u)-1$ 是包含曲线 $\operatorname{im}(u)$ 的射影子空间的最小维数。由于
$r(u)-1\le d$，每条二次曲线（即次数 $2$ 的曲线）都位于一个 $2$-平面中，每条三次曲线
都位于一个 $3$-平面中，依此类推。

用矩阵 $F(u)$ 表示，$G$ 的元素通过矩阵乘法作用：
\[
  F(g\circ u)=gF(u),
  \qquad g\in G.
\]
当 $d\le n$ 时，这个作用在所有满足 $r(u)=d+1$ 的 $u$ 所成的集合上传递。特别地，当
$d=n$ 时，这个轨道由如下\emph{\BookRecordIndex{T0387}有理正规曲线} $C_d$ 生成：
\[
  z\longmapsto[1:z:z^2:\cdots:z^d].
\]
（当 $d=3$ 时，这条曲线也称为\emph{\BookRecordIndex{T0258}扭曲三次曲线}。）
\end{remark}

\begin{exercise}\label{ex:7-4-7}
考虑 $\mathbb{CP}^n$ 中一条满足 $r(u)-1<n$ 的\BookRecordIndex{T0386}有理曲线。证明它位于
$\mathbb{CP}^n$ 的一个 $(r(u)-1)$-维线性子空间 $H$ 中。进一步证明，它是一条具有最大秩
\[
  r=\min(d+1,n+1)
\]
的曲线在适当投影
\sourcefootnote{原书把矩阵 $F(u)$ 的最大秩印作
$\max(d+1,n+1)$；但 $F(u)$ 是 $(n+1)\times(d+1)$ 矩阵，其秩至多为
$\min(d+1,n+1)$，故唯一补正。}
\[
  \pi_Q:\mathbb{CP}^n\setminus Q\longrightarrow H\subset\mathbb{CP}^n
\]
下的像，其中 $Q$ 是 $\mathbb{CP}^n\setminus H$ 中的一个 $(n-r(u))$-维平面，而
$\pi_Q(q)$ 表示 $H$ 与由 $Q$ 及点 $q\notin Q$ 张成的线性子空间的唯一
交点。

\emph{提示：}先考虑 $d=n$ 的情形。
\end{exercise}

\paragraph{$\mathbb{CP}^2$ 中穿过 $3d-1$ 个点的曲线计数。}
本节中强调环境流形很重要，因此我们采用记号
\[
  \mathcal M_{0,k}(M,A;J):=\mathcal M_{0,k}(A;J)
\]
表示 $M$ 中代表类 $A$ 的带标记 $J$-全纯球面\BookRecordIndex{T0246}模空间。考虑\BookRecordIndex{T0283}求值映射
\[
  \operatorname{ev}_{J_0}:
  \mathcal M^*_{0,k}(\mathbb{CP}^n,dL;J_0)
  \longrightarrow(\mathbb{CP}^n)^k.
\]

由于 $J_0$ 可积，\BookRecordIndex{T0246}模空间 $\mathcal M^*_{0,k}(\mathbb{CP}^n,dL;J_0)$ 有自然的复结构
（见注~\ref{rem:3-2-6}），而且\BookRecordIndex{T0283}求值映射 $\operatorname{ev}_{J_0}$ 是全纯的。当 $n=2$ 且
$k=3d-1$ 时，\BookRecordIndex{T0246}模空间的复维数为
\[
  \dim_{\mathbb C}\mathcal M^*_{0,3d-1}(\mathbb{CP}^2,dL;J_0)=6d-2.
\]
注意，这也是 $(\mathbb{CP}^2)^{3d-1}$ 的维数。相应的\BookRecordIndex{T0143}格罗莫夫--威滕不变量
\[
  N_d:=\mathrm{GW}^{\mathbb{CP}^2}_{dL,3d-1}
  (c^2,\ldots,c^2)
\]
计数 $\mathbb{CP}^2$ 中穿过 $3d-1$ 个一般点的次数 $d$ 曲线数。（回忆例
\ref{ex:7-1-14}：$c$ 表示 $H^2(\mathbb{CP}^2;\mathbb Z)$ 的正生成元。）这个数也可解释为
\BookRecordIndex{T0283}求值映射 $\operatorname{ev}_{J_0}$ 的度数；它是良定义的，因为由命题~\ref{prop:7-4-3}
和定理~\ref{thm:6-6-1}，该映射是\BookRecordIndex{T0337}伪循环。全纯映射在\BookRecordIndex{T0405}正则点保持定向，因此 $N_d\ge0$，
并且等号成立当且仅当 $\operatorname{ev}_{J_0}$ 处处不具有最大秩。

\begin{proposition}\label{prop:7-4-8}
对每个整数 $d\ge1$，都有 $N_d>0$。
\sourcefootnote{原书把范围印作 $d\ge0$；但上文
$N_d:=\mathrm{GW}^{\mathbb{CP}^2}_{dL,3d-1}$ 在 $d=0$ 时会有负的\BookRecordIndex{T0066}标记点数，因而
$N_0$ 未由该式定义，且下述证明从正次数开始。故唯一补为 $d\ge1$。}
\end{proposition}

下面说明 Harris--Morrison \srcite{168} 对这一结果的证明。这一论证也可用来展示本书
所讲的全纯曲线辛方法与代数几何方法在观点上的差别。因为我们不假设读者有任何代数几何
知识，所以会解释许多初等细节。

设 $(M,J_0,\omega_0)$ 是 凯勒 曲面。$M$ 中的\emph{\BookRecordIndex{T0389}有效除子}是一个形式和
\[
  C=\sum_i m_iC_i,
\]
其中 $m_i$ 是正整数，每个 $C_i$ 都是 $M$ 中不可约的复一维子簇。用本书的语言，这意味着
\[
  C_i=u_i(\Sigma_i)
\]
是一条单纯 $J_0$-全纯曲线 $u_i:\Sigma_i\to M$ 的像，其中 $\Sigma_i$ 是紧致、连通、无边
Riemann 曲面。在上述 $C$ 的表达式中，我们假设 $i\ne j$ 时 $C_i\ne C_j$。每个\BookRecordIndex{T0389}有效除子
$C$ 都可描述为某个全纯线丛 $L\to M$ 的一个非零全纯截面 $s$ 的零点集；而且复丛同构意义
下的偶对 $(L,s)$ 由 $C$ 唯一确定。因此，每个除子都决定一个上同调类
$c_1(L)\in H^2(M;\mathbb Z)$；用我们的术语说，它是由 $J_0$-全纯曲线的形式和所代表的
同调类
\[
  A=\sum_i m_iA_i
\]
的 Poincaré 对偶。全纯曲线与线丛全纯截面之间的这种对偶，是研究 凯勒 曲面中全纯曲线
的有力工具；在殆复或辛范畴中，没有同等程度的工具可用。

例如，考虑 $M=\mathbb{CP}^2$，令 $L\to\mathbb{CP}^2$ 表示\BookRecordIndex{T0414}重言线丛的第 $d$ 次张量幂。
具体地说，$L$ 是 $\mathbb C^3\times\mathbb C$ 对 $\mathbb C^*$ 取商所得的空间，其中
$\lambda\in\mathbb C^*$ 的作用为
\[
  \lambda\mathbin{\boldsymbol\cdot}(z_0,z_1,z_2,\zeta)
  :=(\lambda z_0,\lambda z_1,\lambda z_2,\lambda^d\zeta).
\]
所以 $L$ 的一个全纯截面是次数为 $d$ 的复系数齐次多项式
\[
  f(z_0,z_1,z_2)
  =\sum_{i+j+k=d}c_{ijk}z_0^iz_1^jz_2^k.
\]
记这种多项式所成的空间为 $\mathcal F_d$。为了描述全纯截面与除子之间的关系，需要使用
代数几何中的三个基本事实；Atiyah--MacDonald \srcite{21} 是很好的参考资料。

\emph{（唯一分解）}多项式环 $\mathbb C[z_0,\ldots,z_n]$ 是\BookRecordIndex{T0332}唯一分解整环。这意味着：若
$f=f_0f_1$ 蕴含 $f_0,f_1$ 中至少一个是常数，就称多项式 $f$ \emph{不可约}。因此，按定义，
每个多项式都\BookRecordIndex{T0213}可分解成不可约因子。唯一分解性质断言，这种分解在重排因子和乘以非零常数的
意义下是唯一的。

\emph{（Noether 性）}多项式环 $\mathbb C[z_0,\ldots,z_n]$ 是 \BookRecordIndex{T0043}Noether 环。这意味着每个
理想都是有限生成的。

\emph{（\BookRecordIndex{T0236}零点定理）}对每个理想 $\mathcal I\subset\mathbb C[z_0,\ldots,z_n]$，记
$\mathcal V(\mathcal I)\subset\mathbb C^{n+1}$ 为 $\mathcal I$ 中元素的公共零点集，并记
$\mathcal I(\mathcal V(\mathcal I))$ 为
$\mathbb C[z_0,\ldots,z_n]$ 中所有在 $\mathcal V(\mathcal I)$ 上消失的多项式所成的理想。
Hilbert \BookRecordIndex{T0236}零点定理断言
\[
  \mathcal I(\mathcal V(\mathcal I))=\sqrt{\mathcal I},
\]
其中\BookRecordIndex{T0146}根理想 $\sqrt{\mathcal I}$ 是所有满足如下性质的多项式
$f\in\mathbb C[z_0,\ldots,z_n]$ 所成的集合：存在正整数 $k$，使 $f^k\in\mathcal I$。

现在回到 $n=2$ 的情形。齐次多项式的零点集在标量乘法下不变，因而可等同于
$\mathbb{CP}^2$ 的子集。其次，齐次多项式的不可约因子也齐次。第三，可以证明，非零不可约
齐次多项式 $f\in\mathcal F_d$ 的零点集
\[
  C_f:=\bigl\{[z_0:z_1:z_2]\in\mathbb{CP}^2\mid
  f(z_0,z_1,z_2)=0\bigr\}
\]
是一条定义在闭连通 Riemann 曲面 $\Sigma$ 上的 $J_0$-全纯曲线
$u:\Sigma\to\mathbb{CP}^2$ 的像（一个简单例子见习题~\ref{ex:7-1-16}）。

把 Hilbert \BookRecordIndex{T0236}零点定理用于由 $f$ 生成的主理想，可知当 $f$ 不可约时，
$\mathbb C^3$ 上每个在 $C_f$ 上消失的多项式都能被 $f$ 整除。因此，次数为 $d$ 的不可约
齐次多项式 $f$ 由其零点集 $C_f$ 在相差常数因子的意义下唯一确定。在可约情形，我们用
$C_f$ 表示 $f$ 的各不可约因子零点集的形式和。这是 $\mathbb{CP}^2$ 中的\BookRecordIndex{T0389}有效除子；由唯一
分解性质和\BookRecordIndex{T0236}零点定理，每个非零多项式 $f\in\mathcal F_d$ 都由其除子 $C_f$ 在相差常数因子的
意义下唯一确定。因此，$\mathbb{CP}^2$ 中次数为 $d$ 的\BookRecordIndex{T0389}有效除子之集可自然地等同于射影空间
\[
  \mathcal P_d:=\mathbb P\mathcal F_d.
\]
注意，$\mathcal P_d$ 的复维数为
\[
  \dim_{\mathbb C}\mathcal P_d=m:=\frac12d(d+3).
\]
$\mathcal P_d$ 在元素 $[f]$ 处的切空间可等同于商空间
$\mathcal F_d/\mathbb C f$。不过，这个等同并非典范的，而依赖于 $f$：把 $f$ 换成
$\lambda f$，会把切向量 $\phi\in\mathcal F_d/\mathbb C f$ 换成 $\lambda\phi$。

若想从\BookRecordIndex{T0389}有效除子的代数几何描述中恢复 $J_0$-全纯球面的\BookRecordIndex{T0246}模空间，就会考虑那些非零多项式
$f\in\mathcal F_d$，其零点集 $C_f$ 是 $\mathbb{CP}^2$ 中代表类 $dL$ 的单纯
$J_0$-全纯球面的像。这恰好是所谓的\emph{\BookRecordIndex{T0049}Severi 簇}
\[
  \mathcal V_d
  :=\bigl\{[f]\in\mathcal P_d\mid
  f\text{ 不可约，且 }C_f\text{ 的亏格为零}\bigr\}
\]
（见 Harris--Morrison \srcite{168}，第 1.F 章）。集合 $\mathcal V_d$ 相当复杂，因而值得
专注于开子集 $\mathcal U_d\subset\mathcal V_d$：它由所有不可约的亏格零次数 $d$ \BookRecordIndex{T0389}有效除子
组成，并且这些除子的奇点全都是简单\BookRecordIndex{T0191}节点。代表元 $f\in\mathcal F_d$ 的零点集 $C_f$ 是一个
带\BookRecordIndex{T0167}横截自交点的浸入球面，每点至多有两个分支。注~\ref{rem:7-4-2} 断言，此时\BookRecordIndex{T0167}横截自交点的
数目为
\begin{equation}\label{eq:7-4-3}
  \delta_d:=\frac12(d-1)(d-2),
\end{equation}
反过来，每条具有 $\delta_d$ 个\BookRecordIndex{T0167}横截自交点的不可约次数 $d$ 的 $J_0$-全纯曲线都必为浸入
球面。因此，$\mathcal U_d$ 有如下初等描述。给定非零齐次多项式
$f\in\mathcal F_d$，记 $S_f\subset C_f$ 为 $C_f$ 的所有简单\BookRecordIndex{T0191}节点之集，即
\[
  S_f:=\left\{[z_0:z_1:z_2]\in\mathbb{CP}^2\ \middle|\
  f(z)=0,\ \frac{\partial f}{\partial z_i}(z)=0,\
  \operatorname{rank}\left(\frac{\partial^2f}{\partial z_i\partial z_j}(z)\right)=2
  \right\}.
\]
（例如，$C_{z_0z_1}$ 在 $[0:0:1]$ 处有简单\BookRecordIndex{T0191}节点，而
$C_{z_0^2-z_1^3}$ 没有。\sourcefootnote{原书在此直接以非齐次式
$z_0^2-z_1^3$ 标记 $\mathbb{CP}^2$ 中的曲线；这可理解为仿射图 $z_2=1$ 上的局部尖点
方程，若写成齐次方程则为 $z_0^2z_2-z_1^3=0$。}）于是
\begin{equation}\label{eq:7-4-4}
  \mathcal U_d
  =\bigl\{[f]\in\mathcal P_d\mid f\text{ 不可约，}\#S_f=\delta_d\bigr\}.
\end{equation}
令 $\mathcal U_d^*\subset\mathcal U_d$ 表示所有如下等价类 $[f]\in\mathcal U_d$ 的子集：
不可约多项式 $f$ 的\BookRecordIndex{T0191}节点集 $S_f$ 处于一般位置，即
\begin{equation}\label{eq:7-4-5}
  \mathcal U_d^*:=\bigl\{[f]\in\mathcal U_d\mid S_f\text{ 正则}\bigr\}.
\end{equation}
这里称有限集 $S\subset\mathbb{CP}^2$ \emph{正则}，如果
$\dim\mathcal F_d/\mathcal H_S=\#S$，其中 $\mathcal H_S\subset\mathcal F_d$ 是线性子空间
\begin{equation}\label{eq:7-4-6}
  \mathcal H_S
  :=\bigl\{\phi\in\mathcal F_d\mid\phi(p)=0\ \forall p\in S\bigr\}.
\end{equation}
给定 $\mathbb C^3$ 上的齐次函数 $\phi$ 和点 $p\in\mathbb{CP}^2$，方程 $\phi(p)=0$ 是指
$\phi$ 在 $p$ 的某个提升、从而每个提升上消失。对每个 $p\in S$ 选取一个提升
$z_p\in\mathbb C^3$ 后，在各点 $z_p$ 求值便给出以 $\mathcal H_S$ 为核的线性映射
$\mathcal F_d\to\mathbb C^S$。所以，$S$ 正则恰好等价于该映射满射。例如，
$\mathbb{CP}^2$ 中四个不同的共线点所成的集合在 $d\ge3$ 时正则，而在 $d\le2$ 时不正则。
这是因为，次数 $d$ 的单变量多项式中有四个指定零点的多项式空间，在所有次数 $d$ 多项式
所成空间中的余维恰好在 $d\ge3$ 时为 $4$。

我们对命题~\ref{prop:7-4-8} 的证明分五步。前两步先观察到 $\mathcal U_d^*$ 是
$\mathcal P_d$ 的复维数 $3d-1$ 的光滑子流形。为了说明 $\mathcal U_d^*$ 非空，我们在步骤
3 和 4 中从 $d$ 条直线的并出发，通过消解 $d-1$ 个\BookRecordIndex{T0191}节点来构造 $\mathcal U_d^*$ 的元素。
这是第~10 章辛\BookRecordIndex{T0254}粘合构造的代数版本。最后，我们用一个运用 Hilbert \BookRecordIndex{T0236}零点定理的\BookRecordIndex{T0166}横截性论证，
证明从 $\mathcal U_d^*$ 出发的\BookRecordIndex{T0283}求值映射具有正的度数。

\begin{proof}[命题~\ref{prop:7-4-8} 的证明]
考虑集合
\[
  \mathcal S_d
  :=\bigl\{([f],p)\in\mathcal P_d\times\mathbb{CP}^2
  \mid f(p)=0,\ df(p)=0\bigr\}.
\]
它由所有偶对 $([f],p)\in\mathcal P_d\times\mathbb{CP}^2$ 组成，其中 $p$ 是 $C_f$ 的
奇点。

\emph{步骤 1. $\mathcal S_d$ 是 $\mathcal P_d\times\mathbb{CP}^2$ 的余维 $3$ 的光滑闭
子流形。此外，投影 $\pi:\mathcal S_d\to\mathcal P_d$ 在
$([f],p)\in\mathcal S_d$ 附近是浸入，当且仅当 $p$ 是 $C_f$ 的简单\BookRecordIndex{T0191}节点。此时，$\pi$ 在
$([f],p)$ 处的微分之像为}
\[
  \operatorname{im}d\pi([f],p)=\mathcal H_p/\mathbb C f,
  \qquad
  \mathcal H_p:=\bigl\{\phi\in\mathcal F_d\mid\phi(p)=0\bigr\}.
\]

$\mathbb{CP}^2$ 上的仿射坐标卡，是一个像不经过原点的仿射映射
$\mathbb C^2\to\mathbb C^3$。在这样的坐标 $(x,y)\in\mathbb C^2$ 中，
$\mathcal F_d$ 里的多项式具有形式
\[
  f_a(x,y)=\sum_{i+j\le d}a_{ij}x^iy^j.
\]
共有 $m+1$ 个系数 $a_{ij}$，其中 $m:=(d^2+3d)/2$。若
$([f_a],x,y)\in\mathcal S_d$，则 $f_a$ 不可能是线性的。因此，可以选取一对指标
$i_0,j_0\ge0$，满足 $1<i_0+j_0\le d$，并施加条件 $a_{i_0j_0}=1$，以此规范化 $f_a$。
这在 $\mathcal P_d\times\mathbb{CP}^2$ 上给出局部坐标
$(a,x,y)\in\mathbb C^m\times\mathbb C^2$。在这些坐标中，点
$([f_a],x,y)$ 属于 $\mathcal S_d$，当且仅当
\[
  f_a(x,y)=0,
  \qquad \partial_xf_a(x,y)=0,
  \qquad \partial_yf_a(x,y)=0.
\]
我们强调，给定任意点 $p\in\mathbb{CP}^2$，都可选取仿射坐标卡，使 $p$ 的坐标为
$x=0,y=0$。

现在考虑映射 $F:\mathbb C^m\times\mathbb C^2\to\mathbb C^3$，定义为
\[
  F(a,x,y):=\bigl(f_a(x,y),\partial_xf_a(x,y),\partial_yf_a(x,y)\bigr).
\]
该映射的零点集就是 $\mathcal S_d$。不失一般性，假设 $F(a,0,0)=0$。于是
$a_{00}=a_{10}=a_{01}=0$，而 $F$ 在 $(a,0,0)$ 处的微分为
\[
  dF(a,0,0)=
  \begin{pmatrix}
    1&0&0&0&\cdots&0&0&0\\
    0&1&0&0&\cdots&0&2a_{20}&a_{11}\\
    0&0&1&0&\cdots&0&a_{11}&2a_{02}
  \end{pmatrix},
\]
其中前三列是关于 $a_{00},a_{10},a_{01}$ 的导数，最后两列是关于 $x,y$ 的导数。因此
$dF(a,0,0)$ 满射，从而 $\mathcal S_d$ 是
$\mathcal P_d\times\mathbb{CP}^2$ 的余维 $3$ 子流形。

仍假设 $F(a,0,0)=0$。容易看出，$C:=f_a^{-1}(0)$ 在原点的奇点是简单\BookRecordIndex{T0191}节点，当且仅当
\[
  a_{11}^2-4a_{20}a_{02}\ne0.
\]
$dF(a,0,0)$ 的公式表明，这又等价于投影
\[
  \ker dF(a,0,0)\longrightarrow\mathbb C^m,
  \qquad(\alpha,\xi,\eta)\longmapsto\alpha
\]
是单射。此外，这个投影的像是所有满足 $\alpha_{00}=0$ 的
$\alpha\in\mathbb C^m$ 所成的集合。用内蕴语言说，这意味着
$([f],p)\in\mathcal S_d$ 在 $p$ 处有简单\BookRecordIndex{T0191}节点，当且仅当投影
$\pi:\mathcal S_d\to\mathcal P_d$ 在 $([f],p)$ 处的微分是单射；在这种情形，
$d\pi([f],p)$ 的像与 $\mathcal P_d$ 中所有在 $p$ 点消失的多项式所成的子流形相切。
步骤 1 得证。

\emph{步骤 2. $\mathcal U_d^*$ 是 $\mathcal P_d$ 中复余维
$\delta:=\delta_d$ 的子流形，而且对每个 $[f]\in\mathcal U_d^*$，}
\[
  T_{[f]}\mathcal U_d^*=\mathcal H_{S_f}/\mathbb C f.
\]

取 $[f]\in\mathcal U_d^*$，并记
$\{p_1,\ldots,p_\delta\}:=S_f\subset C_f$ 为 $C_f$ 的简单\BookRecordIndex{T0191}节点集。于是对
$j=1,\ldots,\delta$，都有 $([f],p_j)\in\mathcal S_d$；由步骤 1，投影
$\pi:\mathcal S_d\to\mathcal P_d$ 在每个 $([f],p_j)$ 处的微分都单射。
\sourcefootnote{原书此句末误印作“for every $k$”；本段只以 $j$ 标号\BookRecordIndex{T0191}节点，故据紧邻定义
唯一补为“对每个 $j$”。}选取 $([f],p_j)$ 的两两不交的开邻域
$\mathcal N_j\subset\mathcal S_d$，使 $\pi$ 在每个 $\mathcal N_j$ 上的限制都是嵌入。
由于 $S_f$ 正则，只要这些邻域取得充分小，这些限制就横截相交。因此，集合
\[
  \mathcal W:=\bigcap_{j=1}^\delta\pi(\mathcal N_j)
\]
是 $\mathcal P_d$ 中复余维为 $\delta$ 的子流形。\sourcefootnote{原书此处误称 $\mathcal W$
的复维数为 $\delta$；它是 $\delta$ 个横截超曲面的交，故复余维为 $\delta$，这也与本步
陈述及下一步的维数计算一致。}现在由注~\ref{rem:7-4-2} 回忆，
$C_f$ 除了这些 $p_j$ 以外没有其他奇点。所以，如果 $f'\in\mathcal F_d$ 接近 $f$ 且同样有
$\delta$ 个简单\BookRecordIndex{T0191}节点，那么这些\BookRecordIndex{T0191}节点必接近各 $p_j$，从而 $[f']$ 必属于 $\mathcal W$。
反过来，按定义 $\mathcal W\subset\mathcal U_d$。因此 $\mathcal W$ 是
$\mathcal U_d^*$ 中 $[f]$ 的一个邻域。这证明了步骤 2。

下面说明如何从 $k+1$ 条曲线的可约并
$C_{f_0}\cup\cdots\cup C_{f_k}$ 构造一条\BookRecordIndex{T0070}不可约曲线：在每个集合
$C_{f_i}\cap C_{f_0}$（$i=1,\ldots,k$）中选一个交点，共消解（即光滑化）$k$ 个交点。

\emph{步骤 3. 设 $f_i\in\mathcal F_{d_i}$ 对 $i=0,\ldots,k$ 都不可约。定义}
\[
  f:=f_0f_1\cdots f_k\in\mathcal F_d,
  \qquad d:=d_0+d_1+\cdots+d_k,
\]
\sourcefootnote{原书把多项式 $f=f_0\cdots f_k$ 印作属于射影化空间 $\mathcal P_d$；
但 $f_i$ 及其乘积都是 $\mathcal F$ 中的多项式，而下文才使用其射影类 $[f]\in\mathcal P_d$，
故唯一补为 $f\in\mathcal F_d$。}
\emph{并假设 $C_f$ 的奇点全都是简单\BookRecordIndex{T0191}节点。对每个 $i=1,\ldots,k$，选取 $f$ 的一个简单
\BookRecordIndex{T0191}节点 $p_i\in C_{f_0}\cap C_{f_i}$。则存在 $([f],p_i)$ 的邻域
$\mathcal N_i\subset\mathcal S_d$ 和 $[f]$ 的邻域 $\mathcal W\subset\mathcal P_d$，使得}
\[
  [f']\in\mathcal W\setminus\bigcup_{i=1}^k\pi(\mathcal N_i)
  \qquad\Longrightarrow\qquad f'\text{ 不可约}.
\]

由步骤 1，可把 $\mathcal N_i$ 取得足够小，使投影
$\pi:\mathcal N_i\to\mathcal P_d$ 是嵌入。对每个 $i=1,\ldots,k$，取 $p_i$ 的开邻域
$U_i\subset\mathbb{CP}^2$ 和 $[f]$ 的开邻域 $\mathcal W\subset\mathcal P_d$，使
\[
  (\mathcal W\times U_i)\cap\mathcal S_d\subset\mathcal N_i.
\]
这意味着，若
$[f']\in\mathcal W\setminus\bigcup_i\pi(\mathcal N_i)$，则对每个 $p'\in U_i$，都有
$([f'],p')\notin\mathcal S_d$，所以 $C_{f'}$ 在 $U_i$ 中没有奇点。由于 $C_f$ 的奇点全是
简单\BookRecordIndex{T0191}节点，缩小 $\mathcal W$ 后可以假设：对每个 $[f']\in\mathcal W$，$C_{f'}$ 的奇点
也全是简单\BookRecordIndex{T0191}节点，而且位于 $C_f$ 中那些未消解的简单\BookRecordIndex{T0191}节点附近。现在，对每个 $i$，在
$U_i$ 的边界上固定两个\BookRecordIndex{T0405}正则点 $q_i\in C_{f_i}$ 和 $q_{i0}\in C_{f_0}$。若 $f'$ 充分接近
$f$ 且 $[f']\notin\pi(\mathcal N_i)$，则在 $C_{f'}$ 中存在附近的点
$q_i',q_{i0}'$，它们可由 $C_{f'}$ 中的一条正则道路连接。容易\BookRecordIndex{T0326}推出，$C_{f'}$ 中任意两个
\BookRecordIndex{T0405}正则点都可由其中的正则道路连接，因而 $f'$ 不可约。因此，可以选取 $\mathcal W$，使步骤
3 的结论成立。

\emph{步骤 4. $\mathcal U_d^*$ 非空。}

在 $\mathbb{CP}^2$ 中选取一般位置的 $d$ 条直线 $L_1,\ldots,L_d$（即两两不同，且没有
三重交点）。对 $i=1,\ldots,d$，选取非零线性多项式 $f_i\in\mathcal F_1$，使
$C_{f_i}=L_i$。那么多项式
\[
  f:=f_1\cdots f_d\in\mathcal F_d
\]
有 $\delta+d-1=\tfrac12d(d-1)$ 个简单\BookRecordIndex{T0191}节点，且没有其他奇点。（回忆由
式~\eqref{eq:7-4-3}，$\delta=\delta_d=\tfrac12(d-1)(d-2)$。）记 $f$ 的简单\BookRecordIndex{T0191}节点集为
\[
  S_f:=\{p_{ij}\mid1\le i<j\le d\},
  \qquad p_{ij}=L_i\cap L_j.
\]
当 $i<j$ 时，函数
\[
  g_{ij}:=\prod_{k\ne i,j}f_k
\]
的零点集是 $\bigcup_{k\ne i,j}L_k$，所以它在 $S_f$ 中除 $p_{ij}$ 外的所有点都消失。
选取二次多项式 $q_{ij}\in\mathcal F_2$，使 $q_{ij}$ 在 $p_{ij}$ 处不消失。那么多项式
$q_{ij}g_{ij}\in\mathcal F_d$ 在商空间
$\mathcal F_d/\mathcal H_{S_f}$ 中给出 $\#S_f$ 个线性无关的元素。因此，集合 $S_f$ 正则。

现在选取 $([f],p_{ij})$ 的小邻域 $\mathcal N_{ij}\subset\mathcal S_d$。由于 $S_f$ 正则，
集合
\[
  \mathcal W':=\bigcap_{2\le i<j\le d}\pi(\mathcal N_{ij}),
  \qquad
  \mathcal W_j'':=\mathcal W'\cap\pi(\mathcal N_{1j}),
  \quad 2\le j\le d,
\]
是 $\mathcal P_d$ 的子流形，并且
\[
  \operatorname{codim}_{\mathbb C}\mathcal W'=\delta=\delta_d,
  \qquad
  \operatorname{codim}_{\mathbb C}\mathcal W_j''=\delta+1.
\]
所以 $\mathcal W'':=\mathcal W_2''\cup\cdots\cup\mathcal W_d''$ 是
$\mathcal W'$ 的真子集。由步骤 3，每个充分接近 $f$ 且满足
$[f']\in\mathcal W'\setminus\mathcal W''$ 的元素 $f'\in\mathcal F_d$ 都不可约。此外，对
每个这样的 $f'$，曲线 $C_{f'}$ 恰有 $\delta$ 个简单\BookRecordIndex{T0191}节点，而且这些简单\BookRecordIndex{T0191}节点所成的集合
正则。由注~\ref{rem:7-4-2}，$C_{f'}$ 是浸入球面，因而
$[f']\in\mathcal U_d^*$。这证明了步骤 4。

\emph{步骤 5. 证明命题。}

取 $[f]\in\mathcal U_d^*$。由\BookRecordIndex{T0236}零点定理，$\mathcal F_d$ 中每个在 $C_f$ 上消失的多项式
都是 $f$ 的标量倍。这可写成
\[
  \bigcap_{p\in C_f}\mathcal H_p=\mathbb C f.
\]
按 $\mathcal U_d^*$ 的定义，集合 $S_f$ 正则，所以所有 $p\in S_f$ 对应的超平面
$\mathcal H_p$ 之交为 $\mathcal H_{S_f}$。其维数为
\[
  \dim\mathcal F_d-\#S_f
  =\frac12d(d+3)+1-\delta
  =\frac12d(d+3)+1-\frac12(d-1)(d-2)
  =3d.
\]
归纳可知，存在有限集 $X\subset C_f\setminus S_f$，使得
\[
  \mathcal H_{S_f}\cap\mathcal H_X=\mathbb C f,
  \qquad \#X=3d-1.
\]
另一方面，由步骤 2，$\mathcal U_d^*$ 是维数 $3d-1$ 的光滑流形，其切空间为
$T_{[f]}\mathcal U_d^*=\mathcal H_{S_f}/\mathbb C f$。所以
$\mathcal U_d^*$ 在 $[f]$ 处横截于射影空间
$\mathbb P\mathcal H_X\subset\mathcal P_d$。因而对每个满足 $\#Y=3d-1$ 且充分接近
$X$ 的有限集 $Y\subset\mathbb{CP}^2$，都有
\[
  \mathcal U_d\cap\mathbb P\mathcal H_Y\ne\varnothing.
\]
所以\BookRecordIndex{T0283}求值映射
\[
  \operatorname{ev}_{J_0}:
  \mathcal M^*_{0,3d-1}(\mathbb{CP}^2,dL;J_0)
  \longrightarrow(\mathbb{CP}^2)^{3d-1}
\]
的像包含 $(\mathbb{CP}^2)^{3d-1}$ 的一个开子集。因此，其像中存在一个\BookRecordIndex{T0409}正则值；再由
命题~\ref{prop:7-4-5}，$N_d>0$。命题~\ref{prop:7-4-8} 得证。
\end{proof}

Kontsevich 在文献~\srcite{215} 中发现了一个优美的递推公式，可用于计算 $N_d$。我们将在
下文命题~\ref{prop:7-5-11} 中介绍其思想。

\paragraph{高维情形。}
我们必须证明，\BookRecordIndex{T0145}格罗莫夫--威滕伪循环
\[
  \operatorname{ev}_{J_0}:
  \mathcal M^*_{0,3d-1}(\mathbb{CP}^n,dL;J_0)
  \longrightarrow(\mathbb{CP}^n)^{3d-1}
\]
（其复维数为 $n(d+1)+4(d-1)$）对每个 $d$ 和每个 $n\ge3$ 都代表非零同调类。
\sourcefootnote{原书此处的\BookRecordIndex{T0145}格罗莫夫--威滕伪循环定义域漏印星号；定理~\ref{thm:7-4-1}
及本节前文的同一\BookRecordIndex{T0283}求值映射均以\BookRecordIndex{T0091}单纯曲线\BookRecordIndex{T0246}模空间
$\mathcal M^*_{0,3d-1}$ 为定义域，故唯一补出。}作为热身，先考虑 $d=2$。我们给出两个
证明，第一个是几何证明，第二个包含更多细节。

\begin{proposition}\label{prop:7-4-9}
\begin{equation}\label{eq:7-4-7}
  \mathrm{GW}^{\mathbb{CP}^n}_{2L,5}(c^2,c^2,c^n,c^n,c^n)=1.
\end{equation}
\end{proposition}

\begin{proof}[第一个证明]
我们先断言，恰有一条二次曲线（即次数为 $2$ 的球面）穿过两个复余维为 $2$ 的一般平面
$X_1,X_2$ 和三个一般点 $p_3,p_4,p_5$。事实上，这三个点确定唯一的 $2$-平面 $P$。由
例~\ref{ex:7-2-9}，$P$ 中恰有一条二次曲线穿过这三个点以及两个交点
$X_1\cap P,X_2\cap P$。由注~\ref{rem:7-4-6}，每条二次曲线都位于一个平面中，所以恰有
一条二次曲线 $u$ 满足给定条件。

还须验证\BookRecordIndex{T0283}求值映射
\[
  \operatorname{ev}:
  \mathcal M^*_{0,5}(\mathbb{CP}^n,2L;J_0)
  \longrightarrow(\mathbb{CP}^n)^5
\]
在 $u$ 处横截地满足这些约束。
\sourcefootnote{原书在命题~\ref{prop:7-4-9} 的两个证明中把该映射的定义域简写成
$\mathcal M(\mathbb{CP}^n,2L;J_0)$；但靶有五个因子，且后文的点为
$[u,\mathbf z]$，故必须使用带五个\BookRecordIndex{T0066}标记点的\BookRecordIndex{T0091}单纯曲线\BookRecordIndex{T0246}模空间
$\mathcal M^*_{0,5}(\mathbb{CP}^n,2L;J_0)$。}在交点 $X_i\cap P$ 处，只须覆盖切于 $P$
的方向，这可由例~\ref{ex:7-2-9} 的论证做到；而在三个点处，只要移动这些点并利用二次
曲线的唯一性，就可得到\BookRecordIndex{T0166}横截性。请读者使用第二个证明中的技巧，把这个论证写得更明确。
\end{proof}

\begin{proof}[第二个证明]
对 $n$ 作归纳。$n=2$ 的情形来自例~\ref{ex:7-2-9}。因此设 $n\ge3$，并归纳假设结论对
$n-1$ 已成立。选取五个射影子空间
\[
  X_1,X_2,X_3,X_4,X_5\subset\mathbb{CP}^{n-1},
\]
使 $X_i$ 在 $i=1,2$ 时维数为 $n-3$，在 $i=3,4,5$ 时维数为零。命题
\ref{prop:7-4-5} 的证明表明，可以把 $X_i$ 选得使其乘积
\[
  \mathbf X:=X_1\times\cdots\times X_5
  \subset(\mathbb{CP}^{n-1})^5
\]
对每棵 $5$-\BookRecordIndex{T0084}带标号树 $T$ 和每个同调类
$B\in H_2(\mathbb{CP}^{n-1};\mathbb Z)$，都横截于\BookRecordIndex{T0283}求值映射
\[
  \operatorname{ev}_{T,J_0}:
  \mathcal M^*_{0,T}(\mathbb{CP}^{n-1},B;J_0)
  \longrightarrow(\mathbb{CP}^{n-1})^5.
\]
因为类 $2L$ 的唯一分解是两条相交直线，所以这意味着：每条穿过三个点
$X_3,X_4,X_5$ 之一并穿过两个 $(n-3)$-平面 $X_1,X_2$ 的直线，都与穿过另两个点的直线
不相交；而每条穿过三个点中任意两个的直线，都不与两个 $(n-3)$-平面相交。由归纳假设，
$\mathbb{CP}^{n-1}$ 中恰有一条穿过 $X_1,\ldots,X_5$ 的二次曲线。

现在定义射影空间 $Y_1,\ldots,Y_5\subset\mathbb{CP}^n$：当 $i=3,4,5$ 时令
$Y_i:=X_i$；当 $i=1,2$ 时，令 $Y_i$ 为包含 $X_i$ 和点
\[
  p:=[0:\cdots:0:1]\in\mathbb{CP}^n\setminus\mathbb{CP}^{n-1}
\]
的唯一 $(n-2)$-平面。我们断言，$\mathbb{CP}^n$ 中每个穿过各 $Y_i$ 的次数二\BookRecordIndex{T0340}稳定映射，
都必是 $\mathbb{CP}^{n-1}$ 中的一条二次曲线。首先，若它是二次曲线，则它是次数二全纯映射
$u=[u_0:\cdots:u_n]$ 的像；它在三个点处与超平面 $u_n=0$ 相交，因而 $u_n$ 必恒为零。
还剩一对相交直线的情形。其中一条直线必包含三个点中的两个，比如 $Y_4,Y_5$；另一条
直线必须穿过 $Y_1,Y_2,Y_3$。根据对循环 $X_i$ 的假设，这条直线不能包含在
$\mathbb{CP}^{n-1}$ 中，因而只在点 $Y_3$ 处与 $\mathbb{CP}^{n-1}$ 相交。但这说明两条直线
并不相交，所以不能组成\BookRecordIndex{T0340}稳定映射。由此证明：$\mathbb{CP}^n$ 中恰有一个穿过各 $Y_i$ 的
次数二\BookRecordIndex{T0340}稳定映射，而且它就是 $\mathbb{CP}^{n-1}$ 中的二次曲线。

还须证明\BookRecordIndex{T0283}求值映射
\[
  \operatorname{ev}:
  \mathcal M^*_{0,5}(\mathbb{CP}^n,2L;J_0)
  \longrightarrow(\mathbb{CP}^n)^5
\]
在满足 $u(z_i)\in Y_i$ 的唯一点
$[u,\mathbf z]\in\mathcal M^*_{0,5}(\mathbb{CP}^n,2L;J_0)$ 处横截于
$\mathbf Y:=Y_1\times\cdots\times Y_5$。为此，注意 $\mathbb{CP}^n$ 中 $u$ 附近的每条
次数二曲线都具有形式 $v=[v_0:\cdots:v_n]$，其中
$[v_0:\cdots:v_{n-1}]:\mathbb{CP}^1\to\mathbb{CP}^{n-1}$ 是
$\mathbb{CP}^{n-1}$ 中接近 $u$ 的二次曲线，而 $v_n$ 是二次多项式。这个多项式由它在三个
\BookRecordIndex{T0066}标记点 $z_3,z_4,z_5$ 处的值唯一确定。这恰好给出了在 $\mathbb{CP}^n$ 中证明对
$\mathbf Y$ 横截所需的额外三个参数。因此，该不变量等于 $1$。命题得证。
\end{proof}

\begin{proof}[定理~\ref{thm:7-4-1} 的证明]
我们对 $n$ 作归纳，证明
\begin{equation}\label{eq:7-4-8}
  \mathrm{GW}^{\mathbb{CP}^n}_{dL,3d-1}
  (\underbrace{c^2,\ldots,c^2}_{2d-2},
   \underbrace{c^n,\ldots,c^n}_{d+1})
  \ge N_d.
\end{equation}
当 $n=2$ 时，这由 $N_d$ 的定义成立。当 $n\ge3$ 时，可作如下说明。由于这个不变量计数
穿过 $d+1$ 个一般点 $p_1,\ldots,p_{d+1}$ 的次数 $d$ 曲线，每条有关曲线都位于这些点张成
的 $d$-平面中。其余约束与这个平面相交于 $2d-2$ 个点。由二维结果，共有 $N_d$ 条这样的
曲线。\sourcefootnote{这段概述按字面只适用于 $d=2$：当 $d>2$ 时，复余维为 $2$ 的
约束与 $\mathbb{CP}^d$ 的一般交仍是复余维为 $2$ 的子空间，而不是点。一般情形应以下文
对 $n$ 的逐步归纳为准。}

具体地，设 $n\ge3$，并归纳假设结论对 $n-1$ 已成立。选取 $3d-1$ 个射影子空间
\[
  X_1,\ldots,X_{3d-1}\subset\mathbb{CP}^{n-1},
\]
使 $X_i$ 在 $i=1,\ldots,2d-2$ 时维数为 $n-3$，在
$i=2d-1,\ldots,3d-1$ 时维数为零。命题~\ref{prop:7-4-5} 的证明表明，可以把 $X_i$ 选得
使其乘积
\[
  \mathbf X:=X_1\times\cdots\times X_{3d-1}
  \subset(\mathbb{CP}^{n-1})^{3d-1}
\]
对每棵树 $T$ 和每个同调类 $B\in H_2(\mathbb{CP}^{n-1};\mathbb Z)$，都横截于\BookRecordIndex{T0283}求值映射
\[
  \operatorname{ev}_{T,J_0}:
  \mathcal M^*_{0,T}(\mathbb{CP}^{n-1},B;J_0)
  \longrightarrow(\mathbb{CP}^{n-1})^{3d-1}.
\]
由归纳假设，存在次数为 $d$ 的全纯球面
$u:\mathbb{CP}^1\to\mathbb{CP}^{n-1}$ 和 $3d-1$ 个互异点
$z_1,\ldots,z_{3d-1}\in\mathbb{CP}^1$，使对每个 $i$ 都有
\[
  u(z_i)\in X_i.
\]
事实上，归纳假设断言这样的元组 $(u,z_1,\ldots,z_{3d-1})$ 至少有 $N_d$ 个。现在令
$p:=[0:\cdots:0:1]\in\mathbb{CP}^n\setminus\mathbb{CP}^{n-1}$，并定义射影子空间
$Y_i\subset\mathbb{CP}^n$ 如下：
\[
  Y_i\cap\mathbb{CP}^{n-1}=X_i,
  \qquad p\in Y_i,
  \qquad \dim_{\mathbb C}Y_i=n-2,
  \qquad i=1,\ldots,2d-2,
\]
以及
\[
  Y_i:=X_i,
  \qquad i=2d-1,\ldots,3d-1.
\]
我们断言，\BookRecordIndex{T0283}求值映射
\[
  \operatorname{ev}_{J_0}:
  \mathcal M^*_{0,3d-1}(\mathbb{CP}^n,dL;J_0)
  \longrightarrow(\mathbb{CP}^n)^{3d-1}
\]
\sourcefootnote{原书在本证明的三处\BookRecordIndex{T0283}求值映射中把定义域下标印作 $k$；但本段固定了
$3d-1$ 个约束，靶为 $(\mathbb{CP}^n)^{3d-1}$，且式~\eqref{eq:7-4-8} 也使用
$3d-1$ 个\BookRecordIndex{T0066}标记点，故三处统一补正为 $3d-1$。}
在点 $[u,\mathbf z]$ 处横截于乘积
$\mathbf Y:=Y_1\times\cdots\times Y_{3d-1}$；这里把
$\mathbb{CP}^{n-1}$ 等同于 $\mathbb{CP}^n$ 中的超平面 $y_n=0$，从而把 $u$ 看成映入
$\mathbb{CP}^n$ 的映射。为说明这一点，把 $u$ 写成
\[
  u([s:t])=[u_0(s,t):\cdots:u_{n-1}(s,t):0],
\]
其中每个 $u_i$ 都是次数为 $d$ 的齐次多项式，而且各 $u_i$ 在
$\mathbb C^2\setminus\{0\}$ 中没有公共零点。于是，$u$ 附近的每条次数 $d$ 全纯曲线都具有
形式
\[
  v(z)=[v_0(z):\cdots:v_n(z)],
\]
其中当 $i\le n-1$ 时 $v_i$ 接近 $u_i$，而 $v_n$ 接近零。由于 $v_n$ 次数为 $d$，它由
$d+1$ 个点 $z_{2d-1},\ldots,z_{3d-1}$ 处的值唯一确定。结合
$\mathbf X$ 横截于\BookRecordIndex{T0283}求值映射
\[
  \operatorname{ev}_{J_0}:
  \mathcal M^*_{0,3d-1}(\mathbb{CP}^{n-1},dL;J_0)
  \longrightarrow(\mathbb{CP}^{n-1})^{3d-1}
\]
这一事实，可知 $\mathbf Y$ 横截于
$\operatorname{ev}_{J_0}:\mathcal M^*_{0,3d-1}(\mathbb{CP}^n,dL;J_0)\longrightarrow
(\mathbb{CP}^n)^{3d-1}$。所以定理~\ref{thm:7-4-1} 来自命题~\ref{prop:7-4-5}。
\end{proof}

\begin{remark}\label{rem:7-4-10}
\textup{(i)} 若愿意使用如下事实，上面的证明可大幅缩短：
$\mathbb{CP}^n$ 中两个互补维数的全纯子簇的每个交点，对交数的贡献都是正的。这样便无须
验证曲线横截地满足约束。因此，可以专门考虑 $d+1$ 个点约束都位于同一个 $2$-平面 $Y$
中的情形。由\BookRecordIndex{T0186}交点正性，次数 $d$ 曲线也必须位于 $Y$ 中。其余约束与 $Y$ 相交于 $2d-2$ 个
点，故二维结果表明，穿过这些约束的曲线有 $N_d$ 条。

\textup{(ii)} 定理~\ref{thm:7-4-1} 证明中的归纳论证通过投影，把 $\mathbb{CP}^n$ 的
\BookRecordIndex{T0143}格罗莫夫--威滕不变量联系到 $\mathbb{CP}^{n-1}$（最终联系到 $\mathbb{CP}^2$）的不变量。
假设我们已经知道，$\mathbb{CP}^{n-1}$ 中穿过约束
\[
  p_1,\ldots,p_\ell,Z_1,\ldots,Z_m
\]
的次数 $d$ 曲线有 $N$ 条，其中 $p_i$ 是点，$Z_i$ 是适当维数的复平面。若
$p_0\in\mathbb{CP}^n$ 是任意一个不在 $\mathbb{CP}^{n-1}$ 上的点，则显然的投影
\[
  \pi:\mathbb{CP}^n\setminus\{p_0\}
  \longrightarrow\mathbb{CP}^{n-1}
\]
\sourcefootnote{原书先定义投影中心为 $p_0$，但在该式及随后两个补集中误写成未定义的
$p$；由投影中心及本段末句再次出现的 $p_0$，唯一补为 $p_0$。}
把 $\mathbb{CP}^n\setminus\{p_0\}$ 等同于 $\mathbb{CP}^{n-1}$ 上一个线丛的全空间。沿
次数 $d$ 映射 $u:\mathbb{CP}^1\to\mathbb{CP}^{n-1}$ 拉回这个线丛，所得线丛次数为
$d$。因此，提升给定 $u$ 的次数 $d$ 映射
$v:\mathbb{CP}^1\to\mathbb{CP}^n$ 所成的空间，可等同于
$\mathbb{CP}^1$ 上一个次数 $d$ 线丛的截面空间。穿过 $d+1$ 个点的这种截面恰有一个。
所以，只要 $\ell>d+1$，$\mathbb{CP}^{n-1}$ 中穿过约束
$p_1,\ldots,p_\ell,Z_1,\ldots,Z_m$ 的曲线 $u$，与
$\mathbb{CP}^n\setminus\{p_0\}$ 中满足如下条件的曲线 $v$ 之间存在双射：$v$ 穿过
$d+1$ 个点 $q_1,\ldots,q_{d+1}$，其中 $\pi(q_i)=p_i$；它还与
\[
  \pi^{-1}(p_j),\quad d+1<j\le\ell,
  \qquad\text{以及}\qquad
  \pi^{-1}(Z_j),\quad1\le j\le m,
\]
相交。因此，$\mathbb{CP}^n$ 中相应的\BookRecordIndex{T0143}格罗莫夫--威滕不变量至少与
$\mathbb{CP}^{n-1}$ 中的一样大。

它可能严格更大，因为我们忽略了穿过 $p_0$ 的曲线。例如，若选取四个一般点
$q_1,\ldots,q_4$ 和四条都穿过点
$p_0\in\mathbb{CP}^3\setminus\mathbb{CP}^2$ 的直线
$\ell_1,\ldots,\ell_4$，那么穿过 $q_1,\ldots,q_4,p_0$ 的三次曲线很可能对不变量
\[
  N_3(4,4):=
  \mathrm{GW}^{\mathbb{CP}^3}_{3L,8}
  (c^2,c^2,c^2,c^2,c^3,c^3,c^3,c^3)
\]
有贡献，却不对应于 $\mathbb{CP}^2$ 中的曲线。事实上，习题~\ref{ex:7-5-17} 将证明
$N_3(4,4)=30>12$。
\end{remark}

\begin{exercise}\label{ex:7-4-11}
证明
\[
  \mathrm{GW}^{\mathbb{CP}^3}_{2L,4}(c^3,c^3,c^3,c^3)=0,
\]
尽管维数条件~\eqref{eq:7-1-2} 成立。

\emph{提示：}使用注~\ref{rem:7-4-6} 中每条二次曲线都位于一个 $2$-平面内这一事实。
\end{exercise}

也存在 $k\ne3d-1$ 的非零\BookRecordIndex{T0143}格罗莫夫--威滕不变量。例如，当 $d=1$ 时，穿过两个点和任意
多个额外超平面的直线数给出非零不变量。再举一个不那么平凡的例子：
$\mathbb{CP}^3$ 中与五个一般点和两条直线相交的三次曲线\BookRecordIndex{T0246}模空间维数为零。为判断是否存在
这样的曲线，可以考虑其中四个点共面的特殊情形。这时，该平面中有一个穿过这四个点的
二次曲线单参数族，并且有一条直线穿过另一个点且与给定的两条异面直线相交。因此，存在
一条可约三次曲线，由一条直线和一条二次曲线的连通并组成，并穿过这七个循环。要证明当
这些循环被扰动到一般位置时这条曲线仍然存在，并非平凡之事。可以使用第~10 章所讨论的
\BookRecordIndex{T0254}粘合技巧，也可以使用下一习题所勾勒的初等方法。

\begin{exercise}\label{ex:7-4-12}
证明
\[
  \mathrm{GW}^{\mathbb{CP}^3}_{3L,7}
  (c^2,c^2,c^3,c^3,c^3,c^3,c^3)>0.
\]

\emph{提示：}固定五个互异点 $x_1,\ldots,x_5\in\mathbb{CP}^3$，使其中任意三个不共线、
任意四个不共面。证明：对 $\mathbb{CP}^1$ 中五个互异点的每个元组
$\mathbf z=(z_1,\ldots,z_5)$，都存在唯一的次数三全纯曲线
$u_{\mathbf z}:\mathbb{CP}^1\to\mathbb{CP}^3$，使得
\[
  u_{\mathbf z}(z_i)=x_i,
  \qquad i=1,\ldots,5.
\]
固定点 $z_1,z_2,z_3$，并证明映射
\[
  (z_4,z_5,z_6,z_7)
  \longmapsto
  \bigl(u_{\mathbf z}(z_6),u_{\mathbf z}(z_7)\bigr)
\]
的微分在某个四元组 $(z_4,z_5,z_6,z_7)$ 处秩为 $4$。\BookRecordIndex{T0326}推出这个映射横截于两条一般直线的
乘积 $X_6\times X_7$，且
$u_{\mathbf z}(z_6)\in X_6$、$u_{\mathbf z}(z_7)\in X_7$。由此证明子流形
\[
  \{(x_1,\ldots,x_5)\}\times X_6\times X_7
  \subset(\mathbb{CP}^3)^7
\]
横截于\BookRecordIndex{T0283}求值映射
\[
  \operatorname{ev}_{J_0}:
  \mathcal M^*_{0,7}(\mathbb{CP}^3,3L;J_0)
  \longrightarrow(\mathbb{CP}^3)^7.
\]
然后使用命题~\ref{prop:7-4-5}。（习题~\ref{ex:7-5-17} 将说明该不变量的值为 $5$。）
\end{exercise}
